% =====================================================================
%  Chapter 08 — EARLY WORK I: the Claim-A proof attempt (S18–S27)
% =====================================================================
\chapter{Early Work I: The Claim-A Proof Attempt}
\label{ch:claim-a}

\begin{record}
\textbf{Classification.} This chapter, and the two that follow, are
\emph{early work}. The material is retained in full in
\texttt{proofs/claim\_a\_3d\_proof\_attempt.tex}, has not been deleted, and is
read forward: Chapter~\ref{ch:inheritance} names what every later session
took from it. Sessions \Sn{18}--\Sn{27}, May 2026.
\end{record}

The programme's numerical work had established, by \Sn{17}, that the
regularity exponent \(\delta\) it cared about was positive in the exact Prize
equations, independent of the regularisation, and independent of resolution
(Chapter~\ref{ch:layer3}). What it had not established was why. This chapter
is the attempt to prove it.

\section{What we assumed}
\label{sec:claima-assumed}

\subsection{Stated hypotheses}

The proof document opens (\texttt{sec:setup}, line~131) with the following, in
force everywhere in this chapter:

\begin{enumerate}[leftmargin=2.2em]
\item The domain is the flat torus \(\T^3=\R^3/(2\pi\Z)^3\), and the equations
  are the exact incompressible Navier--Stokes equations \texttt{eq:NS} with
  constant viscosity \(\nu\) and \(u_0\in C^\infty(\T^3)\). No temperature, no
  variable viscosity: this is Route~2 from the first line.
\item \(u\) is a \emph{Leray--Hopf} weak solution, global in time.
\item \(u\) is a \emph{suitable} weak solution in the sense of
  Caffarelli--Kohn--Nirenberg: the \emph{local energy inequality} holds.
\item The auxiliary scalar \(\theta\) is defined by \texttt{eq:theta}
  (line~150) with \(\theta(\cdot,0)=0\). At \(\varepsilon=0\), \(\theta\) does
  not appear in \texttt{eq:NS}; its equation is an identity satisfied by any
  solution.
\item The target is \emph{Claim~A} (\texttt{def:claimA}, line~163): for some
  \(\delta>0\),
  \[
    \norm{\nu\abs{\nabla u}^2}_{L^{1+\delta}(Q_T)} \le C(\nu,E_0,T)<\infty .
  \]
  The energy inequality gives \(\delta=0\) and nothing more.
\item An equivalence used constantly: Claim~A holds if and only if
  \(\omega\in L^{2+2\delta}(Q_T)\), via Calder\'on--Zygmund, since
  \(\abs{\nabla u}^2\le C\abs{\omega}^2+C\abs{S}^2\) and
  \(\abs{S}^2\le C\abs{\omega}^2\) by incompressibility.
\end{enumerate}

\subsection{Tacit hypotheses, surfaced later}

Four assumptions were in force but not stated. Two were caught inside the
early work itself; two were caught only by the audit of
Chapter~\ref{ch:audit}. They are listed here, at the point where they entered,
rather than being sprung on the reader later.

\begin{center}
\small
\begin{tabular}{@{}p{0.42\textwidth} l p{0.36\textwidth}@{}}
\toprule
Tacit assumption & Surfaced & Consequence \\
\midrule
That an exponent computed by H\"older on \(r^{-1}\iint_{Q(r)}\abs{p}\abs{u}\)
was \(r^{-1/6}\)
  & \Sn{19} & Wrong: the correct optimal exponent is \(r^{-1/2}\), via
    \(L^{5/3}\times L^{5/2}\). Corrected inside \S9. \\[3pt]
That the shear layer's spatial mean vanishes over every ball
  & \Sn{23} & False: \(\int_{B_r}e^{iky}\,dx\neq0\) in 3D. The claim was
    corrected; the shear result was re-routed through direct parabolic
    regularity, which does not need it. \\[3pt]
That a local \emph{enstrophy} inequality is available for Leray--Hopf
solutions
  & \Sn{31} & False. Suitable weak solutions carry a local \emph{energy}
    inequality. Defect~D4. \\[3pt]
That \(\norm{\nabla u}_{L^\infty(B_{\rstar})}\le CM\)
  & \Sn{31} & False. Calder\'on--Zygmund operators are unbounded on
    \(L^\infty\); the correct bound carries a logarithm. Defect~D8 --- and the
    single most consequential line in this book. \\
\bottomrule
\end{tabular}
\end{center}

The two that were caught in-session are worth pausing on, because they show
the early work auditing itself. The \(r^{-1/6}\) error was found by rederiving
the same bound through a different route (the \(\theta\)-Caccioppoli of \S10)
and getting \(r^{-1/2}\); the discrepancy was chased down rather than
averaged. The shear-mean error was found by the \emph{test suite}: a unit test
named \texttt{test\_spatial\_mean\_near\_zero} failed on \(\sin(2\pi y/L)\),
and the mathematics was corrected to match the test rather than the test
relaxed to match the mathematics.

\section{What we considered}
\label{sec:claima-considered}

\subsection{Two routes, and why both were kept}

Two strategies were laid out at the start and pursued in parallel.

\emph{Route~B --- Caccioppoli--Gehring.} Derive a Caccioppoli inequality for
\(\abs{\nabla u}^2\), convert it to a reverse H\"older inequality, and apply
Gehring's lemma to gain a small amount of extra integrability. This is the
standard mechanism by which elliptic and parabolic problems bootstrap
themselves, it is robust, and it produces exactly the \(\delta>0\) Claim~A
asks for. Its cost is that the Caccioppoli inequality for Navier--Stokes
carries a pressure term, and the pressure is nonlocal.

\emph{Route~C --- the \(\theta\)-bootstrap.} Given any \(\delta_0>0\), iterate
a regularity gain on \(\theta\) and show the exponent grows without bound.

Route~C closed immediately, and it is the single cleanest thing in the entire
document. The bootstrap map is
\[
  \mathcal B(\delta) = \frac{1+31\delta}{29-\delta},
\]
and \texttt{thm:routeC} (line~508) proves \(\mathcal B(\delta)>\delta\) for
all \(\delta\in(0,29)\) by the two-line observation that this is equivalent to
\((\delta+1)^2>0\). Hence \(\delta_n\to\infty\) from \emph{any} positive
start, and \texttt{cor:smooth} (line~547) converts that into
\(u\in C^\infty\) by Prodi--Serrin.

That is why the whole chapter is about Route~B. Once Route~C was proved,
the entire Prize reduced --- unconditionally, with no limit and no
regularisation --- to producing one positive number.

\subsection{Three angles on the gap}

Route~B has a single gap, \texttt{gap:small-ball} (line~633): a local Morrey
bound on \(u\) in \(L^{10/3}\) over small balls, with a universal constant.
\Sn{18} attacked it from three directions at once.

\begin{figure}[ht]
\centering
\begin{tikzpicture}[
  >={Latex[length=2mm]},
  ang/.style={draw, rounded corners=2pt, align=center, inner sep=4pt,
              font=\small, text width=36mm, minimum height=13mm},
  gap/.style={draw, very thick, rounded corners=3pt, align=center,
              inner sep=5pt, font=\small\bfseries, fill=panelgrey,
              text width=46mm}
]
\node[ang] (a2) at (-4.6, 2.2)
  {\textbf{Angle 2}\\ CKN \(\varepsilon\)-regularity\\ and capacity};
\node[ang] (a1) at ( 0.0, 3.1)
  {\textbf{Angle 1}\\ CKN local \(A(r)\)\\ iteration};
\node[ang] (a3) at ( 4.6, 2.2)
  {\textbf{Angle 3}\\ \(\theta\)-positivity and\\ the heat kernel};

\node[gap] (G) at (0,0)
  {The single estimate \((\star)\)\\[2pt]
   \(\displaystyle r^{-1}\!\!\iint_{Q(r)}\!\abs{p}\abs{u}\le C\,r^{\alpha}\),
   \(\alpha>0\)};

\draw[->, thick] (a2) -- node[left, font=\scriptsize, align=right, pos=0.55]
  {constant\\ blows up\\ near \(S\)} (G);
\draw[->, thick] (a1) -- node[right, font=\scriptsize, align=left, pos=0.45]
  {error terms \(r^{-1},r^{-2}\)\\ defeat the smallness} (G);
\draw[->, thick] (a3) -- node[right, font=\scriptsize, align=left, pos=0.55]
  {\(\theta\in L^{5/3-}\) only;\\ bootstrap cannot start} (G);

\node[below=5mm of G, font=\scriptsize\itshape, text width=0.72\textwidth,
      align=center, text=rulegrey]
  {\texttt{prop:gap-equivalence} (line~1160): the local Morrey bound, local
   pressure closability, the reverse H\"older inequality and Claim~A are all
   equivalent. Any one of them implies all of them --- which is why three
   independent attacks converge on one estimate.};
\end{tikzpicture}
\caption[The three angles converge]{The three attack angles of \Sn{18} and the
single estimate they all reduce to. Current bound at the time:
\(r^{-1/2}\) (corrected in \Sn{19} from a first-pass \(r^{-1/6}\)); target
\(r^{\alpha}\) with \(\alpha>0\).}
\label{fig:three-angles}
\end{figure}

The convergence was not a coincidence and was proved to be inevitable:
\texttt{prop:gap-equivalence} establishes that four apparently different
conditions are the same condition. On the information available in \Sn{18},
that was encouraging rather than discouraging --- it meant the problem had a
single well-defined centre, and it justified spending the following sessions
attacking that centre from new directions rather than looking for a fourth
angle.

\subsection{Why each subsequent move looked right at the time}

The sequence of choices from \Sn{19} to \Sn{27} has an internal logic, and it
should be recorded as it was, without the benefit of knowing where it ends.

\emph{\Sn{19}: build a pressure-free tool.} If the pressure is the obstacle,
eliminate it. The \(\theta\)-identity \(\nu\abs{\nabla u}^2=\partial_t\theta+
u\cdot\nabla\theta-\nu\Delta\theta\) gives a Caccioppoli inequality with no
pressure on the right-hand side at all (\texttt{prop:theta-cacc}, line~1306).
This was a genuinely new tool, and the reasoning behind it is sound.

\emph{\Sn{20}: find a case where it closes.} The shear layer. A completely
solved special case is worth more than a partial general argument, because it
tells you the mechanism is real.

\emph{\Sn{21}: push outward from the solved case.} Near-shear data by
perturbation.

\emph{\Sn{22}: localise the obstacle.} Divergence-freeness plus
Gagliardo--Nirenberg splits \(u\) on a ball into an oscillating part and its
mean. The oscillating part is \emph{automatically} controlled, with exponent
\(1/2\), better than the \(1/3\) Route~B requires. So the entire difficulty
sits in the spatial mean \(\langle u\rangle_{B_r}\) --- a scalar-valued,
three-component object, and a far smaller target than the field.

\emph{\Sn{23}: the mean obeys an ODE.} Integrating Navier--Stokes over a ball
gives an evolution equation for the mean whose right-hand side is a boundary
integral. Gr\"onwall it.

\emph{\Sn{25}--\Sn{26}: bound the forcing.} Estimate each boundary term.

\emph{\Sn{27}: change direction --- iterate down scales.} Rather than bound a
quantity at each scale, derive a recursion between scales and iterate downward
from the global energy bound. The attraction of this, stated explicitly at the
time, is that it does \emph{not} require a small-data starting condition:
CKN's iteration needs \(A(r_0)<\varepsilon_0\) to begin, whereas a downward
recursion starts from a bound that always holds.

That last move is the one that led directly to Chapter~\ref{ch:prize-claim},
and it is the one the audit judged most harshly. On the information available
in \Sn{27} it was a reasonable idea, and it remains a reasonable idea; the
failure is in the execution, not the conception.

\section{What it produced}
\label{sec:claima-produced}

\subsection{Analytical results}

\begin{table}[ht]
\centering
\footnotesize
\begin{tabular}{@{}l l p{0.40\textwidth} l@{}}
\toprule
Label & Line & Statement & Session \\
\midrule
\texttt{thm:routeC} & 508
  & \(\mathcal B(\delta)=(1+31\delta)/(29-\delta)>\delta\); hence
    \(\delta_n\to\infty\) from any \(\delta_0>0\) & \Sn{17} \\
\texttt{cor:smooth} & 547
  & \(\delta_0>0\Rightarrow\nabla u\in L^p\ \forall p\Rightarrow u\in C^\infty\)
    by Prodi--Serrin & \Sn{17} \\
\texttt{prop:gap-equivalence} & 1160
  & Local Morrey \(\Leftrightarrow\) local pressure \(\Leftrightarrow\) reverse
    H\"older \(\Leftrightarrow\) Claim~A & \Sn{18} \\
\texttt{cor:theta-L53} & --
  & \(\theta\in L^{p}(Q_T)\) for every \(p<5/3\); the endpoint fails
    (the parabolic B\'enilan--Brezis--Crandall barrier) & \Sn{18} \\
\texttt{prop:theta-cacc} & 1306
  & A Caccioppoli inequality for \(\nu\abs{\nabla u}^2\) with \emph{no
    pressure} on the right & \Sn{19} \\
\texttt{lem:shear-heat} & 1611
  & For \(u_0=(f(y),0,0)\), the NS system reduces \emph{exactly} to the 1D
    heat equation & \Sn{20} \\
\texttt{thm:shear-claimA} & 1722
  & \textbf{Claim~A holds for every shear-layer initial datum}
    \(f\in H^1(\T^1)\), for every \(\delta\in(0,1)\) & \Sn{20} \\
\texttt{prop:lambda-zero} & 1779
  & \(\lambda_{\max}=0\) for shear flows: \(\norm{D\Phi_t}\le1+Ct^{1/2}\) & \Sn{20} \\
\texttt{thm:near-shear} & 1961
  & Claim~A with \(\delta_0=1/4\) for \(u_0=(f(y),0,0)+w_0\),
    \(f\in H^2\), \(\norm{w_0}_{H^{1/2}}\le\varepsilon_1\) & \Sn{21} \\
\texttt{lem:pressure-traceless} & 2044
  & \(p=p_T+p_\sigma\) with \(p_T=-R_iR_jT_{ij}\) traceless;
    \(C_{\mathrm{tr}}=\tfrac12C_{\mathrm{CZ}}\) & \Sn{21} \\
\texttt{lem:gn-divfree} & 2211
  & Local Gagliardo--Nirenberg for divergence-free fields:
    \(\norm{u-\langle u\rangle}_{L^{10/3}(B_r)}\le Cr^{2/5}\norm{\nabla u}_{L^2(B_r)}\) & \Sn{22} \\
\texttt{prop:mean-obstacle} & 2268
  & The oscillating part is controlled with exponent \(1/2>1/3\): the whole
    difficulty is the mean & \Sn{22} \\
\texttt{thm:routeB-closure} & 2397
  & Route~B closes \emph{if and only if} the mean-Morrey condition holds & \Sn{22} \\
\texttt{lem:mean-gronwall} & 2580
  & The Gr\"onwall ODE for \(h(t)=\abs{\langle u\rangle_{B_r}}\) & \Sn{23} \\
\texttt{thm:conditional-mean-morrey} & 2657
  & \(\Gamma(r)\le C_1r^\beta\Rightarrow\) mean Morrey with \(\alpha=1-3\beta\)
    \(\Rightarrow\) Claim~A & \Sn{23} \\
\texttt{thm:linfty-closes-gap} & 3490
  & \(u\in L^\infty_{\mathrm{loc}}\Rightarrow\Gamma(r)\le C^*r^{1/2}
    \Rightarrow\) Claim~A \(\Rightarrow u\in C^\infty\) & \Sn{25} \\
\texttt{cor:complete-hierarchy} & 3569
  & The full conditional chain, every arrow proved & \Sn{25} \\
\texttt{lem:pressure-forcing} & 3672
  & The pressure boundary forcing obeys
    \(F_p\le Cr^{-23/15}\norm{p}_{L^{5/3}}\) & \Sn{26} \\
\texttt{prop:scale-recursive} & 3976
  & \(E(r)\le C_1E(2r)^{3/2}+C_2r^{3/10}\), explicit
    \(\alpha=1/2\), \(\beta=3/10\) & \Sn{27} \\
\bottomrule
\end{tabular}
\caption[What the Claim-A attempt produced]{The analytical output of
\Sn{18}--\Sn{27}. Line numbers are current as of \Sn{47}.}
\label{tab:claima-produced}
\end{table}

\subsection{The chain}

By \Sn{25} the early work had assembled a complete conditional chain
(\texttt{cor:complete-hierarchy}, line~3569) in which \emph{every implication
is proved}:

\begin{figure}[ht]
\centering
\begin{tikzpicture}[
  >={Latex[length=2mm]},
  n/.style={draw, rounded corners=2pt, font=\scriptsize, align=center,
            inner sep=3pt, minimum height=8mm, text width=25mm},
  open/.style={n, very thick, fill=panelgrey},
  ar/.style={->, thick}
]
\node[open] (ps) at (0,0) {Prodi--Serrin\\\OPEN};
\node[n] (li) at (3.1,0) {\(u\in L^\infty_{\mathrm{loc}}\)};
\node[n] (ga) at (6.2,0) {\(\Gamma(r)\le Cr^{\beta}\)};
\node[n] (mm) at (9.3,0) {mean Morrey};
\node[n] (rb) at (9.3,-1.7) {Route B closes};
\node[n] (d0) at (6.2,-1.7) {\(\delta_0>0\)};
\node[n] (rc) at (3.1,-1.7) {Route C:\\\(\delta_n\to\infty\)};
\node[n] (sm) at (0,-1.7) {\(u\in C^\infty\)};

\draw[ar] (ps)--(li); \draw[ar] (li)--(ga); \draw[ar] (ga)--(mm);
\draw[ar] (mm)--(rb); \draw[ar] (rb)--(d0); \draw[ar] (d0)--(rc);
\draw[ar] (rc)--(sm);
\draw[ar, dashed] (sm.south) to[out=-90,in=-90, looseness=0.55]
  node[below, font=\scriptsize] {which is the Prize --- and also the
  hypothesis at the far left} (ps.south);
\end{tikzpicture}
\caption[The conditional hierarchy of \Sn{25}]{The conditional hierarchy
assembled by \Sn{25}. Every solid arrow is proved in the document. The dashed
arrow is the point: the chain is a circle. The open step --- does every
Leray--Hopf solution satisfy a Prodi--Serrin condition? --- \emph{is} the Clay
problem, and CKN supplies the chain's entry point only at almost every point,
with a constant that diverges as one approaches the singular set.}
\label{fig:hierarchy-s25}
\end{figure}

That circularity was stated explicitly at the time, not discovered later. The
value of the chain is not that it proves anything; it is that it \emph{relocates}
the difficulty into a condition strictly weaker than Prodi--Serrin ---
\(\Gamma(r)=\int_0^T\bigl(\abs{B_r}^{-1}\!\int_{B_r}\abs{\nabla u}^2\bigr)^{1/2}\dee t\le C_1r^\beta\), a
local \(L^1_tH^{1/2}_{\mathrm{loc}}\)-type statement about \emph{averages} of
the gradient --- and gives a smaller target to aim at.

\subsection{Numerical results from this period}

\begin{center}
\small
\begin{tabular}{@{}l l l l@{}}
\toprule
\texttt{exp\_id} & \texttt{claim\_id} & key metric & verdict \\
\midrule
\texttt{EXP-L4-CZ-TG-032} & --- & \(\varepsilon_{\mathrm{CZ}}=1.00\), Claim~A \(100\%\) & \textsf{PASS} \\
\texttt{EXP-L4-CZ-SHEAR-032} & --- & \(\varepsilon_{\mathrm{CZ}}=1.00\), Claim~A \(100\%\) & \textsf{PASS} \\
\texttt{EXP-L4-MM-TG-002} & --- & \(M_\alpha^{\max}=11.03\), \(\gamma=+0.139\) & \textsf{PASS} \\
\texttt{EXP-L4-MM-SH-002} & --- & \(M_\alpha^{\max}=14.82\), \(\gamma=+0.384\) & \textsf{PASS} \\
\texttt{EXP-L4-GAM-ADV-001} & --- & \(\gamma:0.000\to-3.263\) & \textsf{NOT SUPPRESSED} \\
\bottomrule
\end{tabular}
\end{center}

\medskip
Two remarks. The mean-Morrey scaling exponent \(\gamma\) is positive for both
initial conditions, which is the sign the conditional theorem wants ---
though, as Chapter~\ref{ch:layer4} records, the first pair of runs was
destroyed by sampling bias and the honest magnitudes come from the random-centre
reruns. And the adversarial \(\Gamma\)-search returned a negative verdict which
is \emph{correct as stated and does not bear on the conjecture}, because it
evolves the data under free viscous decay with the nonlinearity removed. That
row still stands unresolved in the database.

Note also, against the house rules of Chapter~\ref{ch:house-rules}: none of
these five rows carries a \texttt{claim\_id}. The tables they live in predate
the column.

\section{What survived and what did not}
\label{sec:claima-survived}

\begin{table}[ht]
\centering
\small
\begin{tabular}{@{}p{0.27\textwidth} p{0.41\textwidth} p{0.24\textwidth}@{}}
\toprule
Object & Content & Status today \\
\midrule
\texttt{thm:routeC}, \texttt{cor:smooth}
  & Bootstrap from any \(\delta_0>0\) to smoothness
  & \PROVED \\[3pt]
\texttt{thm:shear-claimA}, \texttt{cor:shear-smooth}
  & Claim~A and global regularity for shear-layer data
  & \PROVED \\[3pt]
\texttt{prop:lambda-zero}
  & \(\lambda_{\max}=0\) for shear flows
  & \PROVED \\[3pt]
\texttt{thm:near-shear}
  & Claim~A for near-shear data, \(\delta_0=1/4\)
  & \PROVED \\[3pt]
\texttt{prop:theta-cacc}
  & The pressure-free Caccioppoli inequality
  & \PROVED, dormant \\[3pt]
\texttt{lem:pressure-}\newline\texttt{traceless},
  \texttt{prop:traceless-gain}
  & The traceless CZ decomposition, \(C_{\mathrm{tr}}=\tfrac12C_{\mathrm{CZ}}\)
  & \PROVED, dormant \\[3pt]
\texttt{lem:gn-divfree}, \texttt{prop:mean-obstacle},
\texttt{thm:routeB-closure}
  & Route~B closes iff mean Morrey; the oscillating part is free
  & \PROVED, dormant \\[3pt]
\texttt{thm:conditional-}\newline\texttt{mean-morrey},
  \texttt{cor:complete-}\newline\texttt{hierarchy}
  & The full conditional chain
  & \PROVED{} as implications; the chain is a circle \\[3pt]
\texttt{lem:pressure-forcing}, \texttt{thm:forcing-gap}
  & \(F_p\le Cr^{-23/15}\norm{p}_{L^{5/3}}\), better than the naive
    \(r^{-2}\), still divergent
  & \PROVED; does not close the gap \\[3pt]
\texttt{prop:scale-recursive} (\S19)
  & \(E(r)\le C_1E(2r)^{3/2}+C_2r^{3/10}\)
  & \SUPERSEDED{} by \S20, which was itself \WITHDRAWN \\[3pt]
\texttt{thm:no-concentration},\newline\texttt{cor:claim-a-}\newline\texttt{from-cascade}
  & The uniform no-cascade bound and Claim~A from it
  & \CONDITIONAL{\(u\in L^4_{\mathrm{loc}}\)}; superseded \\[3pt]
The three angles of \Sn{18}
  & CKN capacity, \(A(r)\) iteration, heat-kernel positivity
  & All three \OPEN; each has an identified obstruction \\
\bottomrule
\end{tabular}
\caption[Status of the Claim-A early work]{Status today of every substantial
object produced in \Sn{18}--\Sn{27}. Nothing here is withdrawn on its own
account; \S19 is superseded, and the chain of implications stands.}
\label{tab:claima-status}
\end{table}

Three judgements deserve to be stated rather than left implicit in the table.

\emph{The shear result is the programme's strongest unconditional theorem, and
it is small.} \texttt{thm:shear-claimA} genuinely proves Claim~A, and hence
global regularity, for every shear-layer initial datum in \(H^1(\T^1)\). It
does so because the shear layer reduces the full nonlinear system
\emph{exactly} to the one-dimensional heat equation: the advection
\(u_1\partial_xu_1\) vanishes identically, and \(\Delta p=0\) forces \(p\)
constant. That is a complete and correct proof of a case in which the equation
is not really Navier--Stokes any more. The audit of Chapter~\ref{ch:audit}
left it untouched, and it is untouched today.

\emph{Route~C is worth more than it looks.} Two lines of algebra reduce the
entire Prize problem, unconditionally, to producing one positive exponent. It
has never been challenged by any audit and it is still the terminal step any
future closure would use.

\emph{The mean-Morrey and \(\Gamma(r)\) line is not refuted --- it is
dormant.} From \Sn{31} onward the programme worked on vorticity direction
fields and never returned to it. The implications in
\texttt{cor:complete-hierarchy} remain proved; nobody has shown the
\(\Gamma(r)\) condition is unreachable; nobody has reached it. It is
\OPEN{}, and it is the largest body of proved machinery in the document that
no later session uses.

\begin{obsv}[Where this chapter leads]
\label{obs:claima-leads}
The \Sn{27} decision to iterate a recursion \emph{downward} from a bound that
always holds, rather than upward from a smallness hypothesis, is the hinge.
Its gradient-energy version had an exponent of \(3/4\), sublinear, which as
\texttt{rem:prop19-gap} (line~4175) records reduces the argument to the CKN
small-energy regime and gains nothing. The natural next move was to find a
variable whose recursion is superlinear. Two sessions later, the programme
believed it had.
\end{obsv}
